Les aranyes tenen uns òrgans sensibles en els extrems de les potes que els permeten detectar les vibracions que produeixen els insectes que queden atrapats a la seva teranyina. Considereu que en una teranyina el moviment dels insectes és equivalent al que tindrien en un sistema que es mogués amb un moviment harmònic simple (MHS). Hem observat que un insecte de massa 1,58~g atrapat en una teranyina produeix una vibració de 12~Hz.

\figura[0.45]{fig1.png}

\begin{apartat}{1}
Calculeu la constant elàstica d'aquesta teranyina.
\end{apartat}

\begin{apartat}{1}
Determineu la massa d'un insecte que, en quedar atrapat a la teranyina, té un període d'oscil·lació de 0,12~s. Calculeu el valor absolut de l'acceleració màxima de l'insecte, durant el temps en què es mou a la teranyina, si l'amplitud de l'oscil·lació és de 2,0~mm.
\end{apartat}
